% RESEARCH_DRAFT_SHA256: f7c1db5f71908ca4c24dbaeff234b5746968eed2842e6904da08f3e76e4e2746
% Layout and source-reference wrappers expanded for the web renderer.
\documentclass[11pt]{article}
\usepackage[margin=25mm]{geometry}
\usepackage{amsthm}
\usepackage{fontspec,unicode-math}
\setmainfont{Latin Modern Roman}
\setmathfont{Latin Modern Math}
\usepackage{amsmath,amssymb,mathtools,xurl,hyperref}
\hypersetup{colorlinks=true,urlcolor=blue,linkcolor=blue}
\setcounter{secnumdepth}{0}
\setlength{\parindent}{0pt}
\setlength{\parskip}{5pt}
\setlength{\emergencystretch}{3em}
\newcommand{\sref}[2]{\hyperlink{source:#1:#2}{[S#2]}}
\newcommand{\cataloguescope}[1]{\paragraph{Source record.}#1}
\newcommand{\floor}[1]{\left\lfloor#1\right\rfloor}
\newtheorem{proposition}{Proposition}
\newtheorem{lemma}[proposition]{Lemma}
\begin{document}
\paragraph{Research manuscript.} OpenAI \texttt{gpt-6-astra}, reasoning \texttt{ultra}; 11 October 2026.
\paragraph{Preserved source body.} The complete catalogue mathematical text and sources follow with layout and source links expanded, including source qualifications and its historical preparation notes. New research begins at the section headed Research attempt.

% UnsolvedMath ID 20002134; source catalogue code AIM-INFRASTRUCTURE-0034
\section{A versioned robustness certificate for positive task impact}\label{20002134}
{\small UnsolvedMath ID 20002134 · Computer Science}
\subsection{Definitions and mathematical statement}

\paragraph{Shared notation.}Standard number systems are $\mathbb N,\mathbb Z,\mathbb Q,\mathbb R,\mathbb C$. Specialized notation and index ranges follow the cited source and the definitions below.


A learning outcome is a specified capability assessed after an activity. A comparison requires an outcome instrument, a time window, and a comparison condition. Improvement in observed scores is distinct from a causal effect of the task, because learners and teaching conditions can differ.
\paragraph{Record scope and qualifications.}
The catalogue title is retained as metadata. Named certificates, contracts, bounds or frameworks in that title are not asserted as results of the original source extract.
No task, instrument, population, or comparison design is given. A causal or statistical model would require additional assumptions.

Catalogue screening limitation: Educational assessment, adoption, or study-design discussion; the supplied question does not state a mathematical target with defined assumptions and outcome.

The source snapshot labels this record partially solved; this is catalogue metadata rather than independent verification of a mathematical result or the current completion of a project.
\paragraph{Statement recorded in the supplied source.}
37. How to tell if a particular task (positively) impacts learning \href{http://aimath.org/pastworkshops/oerassessmentproblems.pdf}{[S2]}
\paragraph{Source status.}
The supplied catalogue labels this source record partially solved; that metadata was checked against the supplied snapshot on 2026-10-11. The label is preserved for provenance and is not an independent verification of present open status or mathematical correctness.
\subsection{Short English statement}
Determine whether a particular learning task improves students' learning.
\subsection{Sources}
\begin{itemize}
\item[S1] Ulam AI and the UnsolvedMath Contributors, UnsolvedMath: A Curated Collection of Open Mathematics Problems, record 20002134 (AIM-INFRASTRUCTURE-0034). Catalogue curation and metadata: CC BY 4.0. The source extract is retained; definitions, English summaries and qualifications are editorial additions.
\url{https://huggingface.co/datasets/ulamai/UnsolvedMath}.
\item[S2] AIM Workshop Problem Lists, source linked by catalogue record AIM-INFRASTRUCTURE-0034.
\url{http://aimath.org/pastworkshops/oerassessmentproblems.pdf}.
\end{itemize}
\paragraph{Preparation and review.}
The original source and UnsolvedMath contributors retain attribution for the source extract. Codex (GPT-6; exact serving revision not exposed) assisted with definitions, English summaries, scope notes and formatting. No new human mathematical authorship or independently verified solution is asserted. Permanent source ID 20002134 and code AIM-INFRASTRUCTURE-0034 checked against the complete supplied registry and current repository. Distinct upstream records, including repeated questions and broad or damaged extracts, are retained for complete coverage of the source CS category. A repeated record is not a new independent conjecture. Complete mathematical-equivalence and current-literature review has not been performed. Preparation and final review follow contribution prompts C and D, with source gaps disclosed. Structural and rendering checks do not establish a new theorem or open status.
\label{end:20002134}

\section{Research attempt}
\subsection{Literature, source context and scope}
The primary AIM notes describe a brainstorming session. The original excerpt and all catalogue qualifications remain above. They do not specify a mathematical theorem or a complete empirical study. The formal model below is an explicitly limited research route introduced for this attempt, not a reconstruction of a missing original conjecture. The supplied catalogue's earlier research summary was inspected as context; its named framework is not adopted as an established solution. The elementary result below is not claimed novel, and no educational effectiveness data were collected.

\subsection{Attempt: the resolution of a small randomized task evaluation}
To test whether a very small randomized comparison can provide arbitrarily strong evidence, consider $N$ clusters and an integer $n$ satisfying $1\le n<N$. Exactly $n$ clusters are assigned to a task and the rest to its comparator, uniformly over the $K=\binom Nn$ assignments. Under Fisher's sharp null, every cluster's outcome is fixed and unaffected by assignment. The statistic is the difference of treated and control means.
\begin{proposition}
The exact nonrandomized one-sided permutation $p$-value is at least $1/K$. When $N=2n$, the two-sided value based on the absolute mean difference is at least $2/K$.
\end{proposition}
\begin{proof}
Under the sharp null the complete outcome vector is the same for all assignments, and each assignment has probability $1/K$. The one-sided tail defining the $p$-value includes the observed assignment itself, so its count is at least one. For equal group sizes, complementing an assignment swaps the groups and changes the statistic's sign while preserving its absolute value. An assignment and its complement are distinct when both groups are nonempty. The absolute-value tail therefore contains at least that pair and has probability at least $2/K$.
\end{proof}
\subsection{Verification and remaining gaps}
For four clusters with two per arm, the smallest possible one-sided value is $1/6$ and the two-sided value is $1/3$; outcomes $(0,0,1,1)$ with the two ones treated attain both. With six clusters split three and three, outcomes $(0,0,0,1,1,1)$ attain $1/20=0.05$ and $1/10=0.1$. Enumeration of all assignments verifies the tails. Thus an extreme positive observed contrast does not circumvent the resolution of a specified small exact test.

This result concerns a sharp no-effect null and this nonrandomized exact test, not every test or a null of zero average effect under heterogeneous effects. It does not establish task efficacy or prescribe a study size. Outcomes, clustering, interference and the meaningful educational contrast remain unspecified by the source.

\section{References and source attribution}
American Institute of Mathematics, workshop discussion notes on open educational resources and assessment, \url{https://aimath.org/pastworkshops/oerassessmentproblems.pdf}. The 16-page primary PDF was opened on 11 October 2026 and its numbered brainstorming item was checked against the supplied record. The source supports the provenance and agenda, not the additional modeling assumptions or an empirical educational conclusion. Catalogue and contributor attribution are retained in the reproduced Sources section. The worked mathematics is proved here without reliance on an unverified external theorem.


\paragraph{Integration review.} Codex (GPT-6; exact serving revision not exposed) reviewed the written arguments and expanded layout and source-reference wrappers. Repository metadata and the production web renderer were checked using Python 3.13 and JavaScript with MathJax 3.2.2. This model-assisted review is not independent human proof verification; no PDF compilation or formal proof checking is asserted.

\section{Completion Estimate}
\noindent\textbf{COMPLETION: 0\%}
No substantiated new progress toward the entire original target is claimed. The worked results and failed route are delimited above; a known result, illustrative model, or auxiliary obstruction is not counted as original-target completion.

\section{Final status and attribution}
\textbf{Attempt status: unresolved.} This describes this research attempt, not a new ruling on the source catalogue's status or on external literature claims.
The mathematics and this self-audit were generated by OpenAI \texttt{gpt-6-astra}, reasoning \texttt{ultra}, on 11 October 2026. Codex (GPT-6; exact serving revision not exposed) orchestrated the batch and contributed formatting/repository review. No human mathematical authorship, independent human proof checking, or formal verification is claimed. The historical preparation paragraphs in the reproduced catalogue body concern that body only.

Contribution prompt A was applied to the complete definition and supplied source record, followed by the adversarial checks of prompt B. Local reproduction files are \texttt{batch008/checks.py}, \texttt{computations.json}, and the per-ID audit JSON. They were not public links at preparation time. Numerical checks are finite and do not prove unrestricted results. No LaTeX engine was found locally; compilation was not performed by this research agent.

\end{document}
